% Generated by Listas/make_statements.py from pratica-lecture-06.md, item 28. Do not edit.
\textbf{Source.} Milgrom \& Roberts (1992), \emph{Economics, Organization and Management}, Ch. 8 Exercises, Mathematical Exercise 1, p.~285 {[}PDF 300{]}. Reproduced verbatim (dollar amounts written with an escaped sign, \texttt{\textbackslash{}\$}). Answer in the source: none in the book.



\begin{quote}

\begin{enumerate}

\def\labelenumi{\arabic{enumi}.}

\tightlist

\item

  In the context of the efficiency wage model, suppose that \(N\) is given by \((1 - \delta^t)/(1 - \delta)\), where \(\delta\) is a discount factor giving the value today of a dollar to be received a year from now and \(t\) is the number of years the worker anticipates remaining with the firm. Thus, \(N\) is a measure of the horizon over which the employee is concerned, appropriately discounted. Suppose the worker expects to be employed for 5 years, that is, \(t = 5\), and evaluates future income using \(\delta = 0.9\), so that \(N\) is approximately equal to 4.1. The worker's opportunity wage (\(\overline{w}\)) is \$20,000 per year, and he or she can make an immediate gain of \$3,000 by leaking the firm's trade secrets to a competitor. The firm can monitor the worker to prevent this, but doing so is costly: raising the probability \(p\) of catching misbehavior by an amount \(\Delta p\) costs \$1,500 \(\Delta p\).

\end{enumerate}



Formulate the firm's problem in selecting the monitoring intensity \(p\) and the wage \(w\) and solve for the optimal values.



Suppose now that the threat of a takeover or bankruptcy endangers the employee's future with the firm, so that the employee now expects to be employed only one more year. Will the contract you found above still deter cheating? Why or why not?



More generally, show that in the efficiency wage model, a lengthened horizon allows providing incentives with a lower wage. That is, show that the optimal wage level \(w^*(N)\) is a nonincreasing function of \(N\) provided the marginal cost of monitoring is positive and increasing in \(p\), that is, provided \(M'(p) > 0\) and \(M''(p) > 0\) for all \(p\).

\end{quote}



\textbf{Further practice (not in the tiers).} Source lines only; the statements are in the previous

Lecture 4 list (\texttt{git\ show\ fe47ead:Listas/pratica-lecture-04.md}), where they were verified.

- Muñoz-Garcia Ch. 4 Ex. 16 (pp.~198--199, ``Two-Period Bilateral Bargaining Game'', with Fig. 4.35) and Ex. 17 (pp.~199--200, ``Alternating Bargaining with a Twist''): costly offers and a shrinking surplus, no discounting.

- Tadelis 11.4 (p.~236, ``Constant Delay Cost'') and 11.6 (p.~237, ``Asymmetric Patience 2''): delay costs instead of discounting; Rubinstein with \(\delta_1\neq\delta_2\).

- MWG 9.B.12 (p.~303) and 9.B.13 with 9.B.7 (pp.~302--303): the infinite-horizon game with \(\delta_1\neq\delta_2\); a cost per offer instead of discounting.

- Muñoz-Garcia Ch. 6 Ex. 8 (p.~303, ``Infinite-Horizon Bargaining Game Between Three Players'').



\textbf{Optional extras (collusion, reputation, wages).}

- Tadelis 10.1 (p.~215, discount cards and tacit collusion, discursive), 10.7 (pp.~217--218, the bartender's quality reputation, a twin of item 18) and 10.11 (p.~219, the brand-name equilibrium of §10.5.3).

- Cabral 9.2 (p.~242, the largest Bertrand cartel at \(\delta = .8\)), 9.5 (p.~242, are price wars irrational?), 9.18 (p.~245, multimarket contact again) and 9.20 (p.~246, observed demand fluctuations: Cabral's version of item 23).

- Cabral 9.17 (p.~245): the book asks to show that the minimum discount factor \(\bar\delta\) is \emph{decreasing} in \(n\); with \(\bar\delta = 1 - 1/n\) (Cabral p.~227) it is increasing. Work it as ``show how \(\bar\delta\) moves with \(n\)''.

- Milgrom \& Roberts Mathematical Exercise 2 (pp.~285--286, voluntary quits in the efficiency-wage model) and Food for Thought 1 (p.~284, wage surveys).

- Gibbons 2.3 (p.~131, Rubinstein with \(\delta_1\neq\delta_2\)) and 2.16 (p.~136, a firm and a union in an infinitely repeated wage-employment game).



\begin{center}\rule{0.5\linewidth}{0.5pt}\end{center}

